\begin{figure}[H]
\centering
\begin{tikzpicture}[
  event/.style={draw=orange!70!black, fill=orange!8, rounded corners=3pt,
    text width=2.55cm, minimum height=1.45cm, align=center, font=\scriptsize,
    inner sep=4pt},
  fix/.style={draw=green!55!black, fill=green!7, rounded corners=3pt,
    text width=2.55cm, minimum height=1.45cm, align=center, font=\scriptsize,
    inner sep=4pt},
  arr/.style={-{Stealth[length=2.4mm,width=1.8mm]}, line width=0.8pt, black!45,
    shorten >=1pt, shorten <=1pt},
]

%% row 1 --- building the dataset
\node[event] (a1) at (-5.55, 0)     {Generation run against an external API};
\node[event] (a2) at (-1.85, 0)     {Provider stopped accepting requests};
\node[event] (a3) at ( 1.85, 0)     {Results written only at the end, so the run was lost};
\node[fix]   (a4) at ( 5.55, 0)     {Retry with backoff, and results saved as they are produced};

%% row 2 --- repairing the dataset, then choosing the training method
\node[event] (b1) at (-5.55, -2.20) {One group of Zalo messages written as narration};
\node[fix]   (b2) at (-1.85, -2.20) {Rewritten offline from the stored scenarios};
\node[event] (b3) at ( 1.85, -2.20) {Ordinary LoRA left 9~MiB free and pointed at 18 hours};
\node[fix]   (b4) at ( 5.55, -2.20) {QLoRA fits the card with room to spare};

%% row 3 --- the training run itself
\node[event] (c1) at (-5.55, -4.40) {Output folder name left out the run ID};
\node[fix]   (c2) at (-1.85, -4.40) {Stopped on purpose and restarted from step zero};
\node[fix]   (c3) at ( 1.85, -4.40) {Both models trained to completion};

\draw[arr] (a1)--(a2);  \draw[arr] (a2)--(a3);  \draw[arr] (a3)--(a4);
\draw[arr] (a4.south) -- ++(0,-0.55) -| (b1.north);
\draw[arr] (b1)--(b2);  \draw[arr] (b2)--(b3);  \draw[arr] (b3)--(b4);
\draw[arr] (b4.south) -- ++(0,-0.55) -| (c1.north);
\draw[arr] (c1)--(c2);  \draw[arr] (c2)--(c3);

\end{tikzpicture}
\caption{The order in which things went wrong, and what was changed after each one. Orange is a failure, green is the fix that followed it. Each row is one stage of the work: building the dataset, repairing it and choosing a training method, and the training run itself.}
\label{fig:failure-recovery-timeline}
\end{figure}
